%======================================================================
\section{Discussion and limitations}
\label{sec:discussion}
%======================================================================

\para{What transfer requires.} Four conditions emerged, and each was
necessary in at least one system. First, the physics must be extensive
at the scales involved. This can be checked before any model is trained,
from ground truth at two sizes: the 2D KS density is $82.0$ and $84.0$ per
base box at $44^2$ and $66^2$, constant within the spread of the truth runs,
while in Kolmogorov flow the truth's density per box rises with the
channel width at every drag ($15.9$, $15.4$, $17.8$ at $\alpha=0.1$; $21.5$, $24.3$,
$25.1$ at $\alpha=0.05$; $25.2$, $25.8$, $27.6$ at $\alpha=0.025$). A local model is
extensive by construction (Remark~\ref{rem:extensive_model}), so wherever
the truth is not, a transfer error is expected and is not a failure of
training. Second, the training data must contain the local statistics of
a large domain. The same architecture trained on a small periodic box
misses $\DKY$ by $15$--$38\%$ in three of the four larger configurations we
evaluated (it happens to be accurate at $\alpha=0.1$, $2\times$), and trained on
windows of one large simulation it transfers (Sec.~\ref{sec:box_vs_window}); the window
training is a limited-area problem, and Proposition~\ref{prop:halo} shows
that the forced halo trains exactly that problem. Third, the learned map
must remain chaotic, and chaotic at the right rate, when it runs freely.
Short rollouts leave room for spurious regimes
(Sec.~\ref{sec:fail_regimes}); long squared-error rollouts damp the growth
rates (Sec.~\ref{sec:fail_mse}); free-rollout statistics compared with the
training data detect the first problem without any truth, and the
truth-free selection rule of Sec.~\ref{sec:canonical_results} chose a
model that is within $0.9$--$5.9\%$ of the truth at every width. Fourth, the
model's reach must cover the interaction range of the dynamics.

\para{The interaction range is the design length.} The light cone set by
the speed at which perturbations spread is far too short to constrain the
reach of a model at these time steps (Sec.~\ref{sec:receptive}). The length
that matters is the distance beyond which neighbouring values stop
improving a one-step prediction: three to four sites, $19$--$25$ length
units, for the 1D latent dynamics, against a trained reach of six. In
Kolmogorov flow most of the unstable directions lie in the band around
the forcing wavenumber, $31.5$ of $36$ at $|\bm k|\in[2,8)$ for $\alpha=0.05$ at
$3\times$, whose wavelengths span one to four sites. Models whose encoder
reaches two sites and whose propagator reaches one fail even when trained
on windows of the very domain they are evaluated on, with errors of
$66$--$100\%$ in $\DKY$ (Appendix~\ref{app:tables}). The rule we take from this
is to measure the interaction range on data, for example by fitting a
shared stencil of increasing width, and to give the encoder and
propagator a reach that covers it with a margin. The reach then sets the
smallest domain the model can be used on, which is what fails at $L=25$.

\para{The parameters describe the law, not the state.} Because every
operator is shared by all sites, the trained parameters encode one local
rule, and their number is fixed by that rule rather than by the domain.
The 2D models have about $135{,}000$ parameters. At $66^2$ such a model
reproduces a Kaplan--Yorke dimension of $188$ to within $0.6\%$ of the truth
(Table~\ref{tab:transfer2d}), and extrapolating the measured density of
$0.043$ per unit area, the attractor at $440^2$, on which the same model is used
for data assimilation in Part~II, has a dimension of roughly $8{,}000$. A
model whose size is set by the state would have to grow with that number;
this one does not. The same weight sharing makes training cheap in data:
every window, and every site within it, is a training example of the same
rule, and the 2D KS models are trained on five trajectories of $25$ time
units. Neither a small absolute parameter count nor size independence is
new by itself: convolutional surrogates and neural operators share weights
too \citep{barsinai2019,kochkov2021}. What the dimension comparison adds is
evidence that the shared rule is the right one, in the sense that its
extension to a larger domain carries the right density of chaos. This is
the inductive bias that extensive chaos calls for: a system made of many
copies of one local process is described most economically by learning
that process once.

\para{Transfer error as a probe of the physics.} Because a local model
is extensive by construction, the difference between its dimension
density and the truth's on a larger domain measures how far the truth is
from extensive at that size, once the model is accurate at the native
size. The weakest drag, where the native-size models scatter most and the
transferred models fall below a rising truth, is where the inverse
cascade is least arrested, and we expect non-local effects to be largest
there. Using
transfer error to locate the scales at which a system stops being
extensive is a question in its own right, which we leave to future work.

\para{What locality buys in data assimilation.} The latent field makes
localization possible in a learned space, and with a per-patch decoder it
makes it exact (Proposition~\ref{prop:blockdiag}). The ensemble a localized
filter needs is then set by the local problem, about 32 members per tiling
at every size, while an unlocalized filter needs more members than the
attractor has unstable directions, a number that grows with the domain
(Secs.~\ref{sec:da_2d} and \ref{sec:da_1d}). The price of exact locality is
an error tied to the patch lattice, which the decoder does not remove; the
shifted tilings remove most of it, because the tilings' errors are only
weakly correlated, and they do so without retraining. The remaining
analysis error is mostly the filter's, not the decoder's: the blended
decoder floor is less than half of it (Sec.~\ref{sec:da_accuracy}). Better
use of the ensemble inside each tiling, and models of the propagator's
error in place of isotropic noise, are the natural next steps.

\para{When to prefer a learned model.} A filter built on the correct
equations is four times more accurate than the latent filter, and one built
on equations with a small or correctable error is still better. The
comparison of Sec.~\ref{sec:da_model_error} locates the boundary: what
matters is the part of the physical model's one-cycle forecast error that
bias correction cannot remove, and the latent filter wins once it exceeds
about a sixth of the observation error. Two features of the comparison
carry over to real systems. The errors that parameter estimation and bias
correction cannot reach are the structural ones, a wrong form of an
operator or a process that is missing or unobserved, and those are the
errors of the models of convection and cloud processes that motivate this
work. And the learned model needs no tuning for the error of the model it
replaces: one setting served every case, while the physics-based filter
worked only in a narrow window of its inflation. The model learned from
noisy fields was slightly better than one learned from clean fields, so
measurement error, the part of the error budget that is usually known, can
be handled in training. A process with memory that the current field does
not reveal limits both filters (Sec.~\ref{sec:da_model_error}); propagators
with access to past states are the obvious remedy on the learned side.

\para{Limitations.} We state them plainly.
\begin{itemize}
\item \emph{Ground truth.} The Kolmogorov truths at $\alpha=0.1$, $1\times$, at
  $\alpha=0.05$, $3\times$, and at all three widths for $\alpha=0.025$ rest on one
  40-unit run each; such runs deviated from their 120-unit continuations
  by $-17\%$ to $+23\%$, and were high in five of seven cases. The 2D KS truths
  average over 40 time units and their running estimates still move by
  $3$--$5\%$; there is no truth at $88^2$.
\item \emph{Kolmogorov geometry.} The channel is widened in $x$ only,
  because the forcing fixes the height; transfer in $y$ is untested.
\item \emph{Models.} Most 2D models have one training seed (the 1D model
  has three, the $\alpha=0.05$ canonical model two); the 1D and 2D
  canonical architectures differ (full reach in 1D, per-patch in 2D); the
  weakest drag remains $8$--$25\%$ off.
\item \emph{What the spectra compare.} The model's Lyapunov spectrum is
  that of the latent map. Comparing it with the PDE's assumes that the
  decoder is a smooth, injective map on the attractor, which the agreement
  of the spectra supports but does not prove.
\item \emph{Reconstruction.} Local neural-field latents have been
  reported to fail on KS at the decoder \citep{aroma2024}. Our autoencoders
  reconstruct the field to $0.01\%$ (1D), $2.6$--$2.9\%$ (2D KS) and $1.2\%$
  (Kolmogorov) of its variance, and the per-patch decoders leave seams
  that appear as excess energy above the patch wavenumber.
\item \emph{Data assimilation.} The experiments are twin experiments in
  one 2D system, with regular, dense observations (four per site) and
  additive Gaussian model noise. With one observation per site the filter
  is still four times better than interpolation and stays bounded over 300
  cycles, but longer runs with sparser networks expose the stability of the
  learned propagator, which we leave to a follow-up study. The model errors
  of the physics-based comparator were chosen, not calibrated to a real
  model, and the learned model was trained on complete, noisy fields of the
  true system, which real data rarely provide. Kolmogorov flow was not
  assimilated.
\end{itemize}

\para{Outlook.} The ingredients carry over to any system whose chaos is
extensive and whose dynamics are local at the scale of a few patches. The
target we have in mind is the assimilation of convection-permitting
simulations of organized convection, such as squall lines, where the
equations of the cloud processes are uncertain and much of the state is
unobserved; other candidates are
convection in large aspect-ratio cells, limited-area atmospheric models
nested in a global one, where the forced halo is already the operating
mode, and high-resolution models whose domain is chosen by the
application rather than by the training data.
% TODO(G-pass): add verified citations for the outlook systems if they are kept.
